% 本文件由 assemble.py 按 _work/plans/ch04.json 逐字组装，请勿手改（改 plan 或 blocks 后重新生成）。
\chapter{曲线运动}
\label{ch:04}

\section{竖直面圆周·绳球}
% ---- block 010-03 (原稿第10页) kind=method ----
\textbf{绳球}\quad $T\perp V$，$W=0$

上到最高点\ $V\ge\sqrt{gr}$

\notefig[0.25]{figures/p010_a.png}

上不去\ $\left\{\begin{array}{ll}\text{上半圆} & \text{斜上抛}\\ \text{下半圆} & \text{单摆}\end{array}\right.$

类绳模型$\to$弹力只能向内

最高点最低点拉力差\ $6mg$


\section{竖直面圆周·杆球}
% ---- block 010-04 (原稿第10页) kind=method ----
\textbf{杆球}\quad 杆力\ $W_{\text{总}}=0$

上到最高点\ $V_{\text{高}}\ge0$

上不去\quad 原路反回 % “反回”疑为“返回”

类杆模型\quad 小球受弹力可内可外

\notefig[0.6]{figures/p010_b.png}
% 图 p010_b：一条大波浪形轨道，沿途分段标注“杆”“绳”“杆”

\notefig[0.35]{figures/p010_c.png}
% 图 p010_c：四个小示意图（圆环形“杆”、拱形“杆”、碗形“绳”、圆形“杆”）


\section{关联速度}
% ---- block 041-01 (原稿第41页) kind=method ----
\textbf{运动关联}% 页面左上角标题(带下划线框线)；页左侧另有一条红色大括号，自本块“人 V”图起括至“同轴同ω”图止

\noindent $V_{\text{船}}\cos\alpha=V_{\text{人}}$\quad 船\unclear{比}人快% 页面右上角两行(第二行中间一字辨认不清)

\begin{itemize}
\item 人匀速船加速
\item 人：减速船才有可能匀速
\end{itemize}

\notefig[0.28]{figures/p041_a.png}
% 图内标注：人 $V$(岸上，水平绳段)；绳方向分速度 $V_{\parallel}$ 改变绳长；船速 $V_1$、夹角 $\alpha$；垂直分速度 $V_{\perp}$(斜向下箭头)
\noindent $V_{\perp}$ 绳子摆动(圆周运动)% 图正下方的标注，因紧邻右侧文字未含在图内

\noindent $V=V_1\cos\alpha$\qquad 只能分解实际速度

\notefig[0.25]{figures/p041_b.png}
% 图内标注：船受力图——$F$(沿绳斜向上，带虚线分量)、$mg$ 向下
\noindent $F\sin\alpha=mg$

% ---- block 041-02 (原稿第41页) kind=method ----
\notefig[0.28]{figures/p041_c.png}
% 图内标注：绳绕过上方定滑轮连接左右两物体；左物体速度 $V_1$(水平向右)，绳与水平夹角 $\alpha$；右物体速度 $V_2$(水平向右)，绳与水平夹角 $\beta$；虚线箭头、沿绳/垂直绳的分量箭头
\[
V_1\cos\alpha=V_2\cos\beta
\]

% ---- block 041-03 (原稿第41页) kind=method ----
\notefig[0.25]{figures/p041_d.png}
% 图内标注：水平面上物体 $B$ 向左运动(标 $V\cos\alpha$ 的分量)，绳绕过上方定滑轮另一端竖直挂物体 $A$(向上箭头)，滑轮处标 $V_A$，夹角 $\alpha$
\noindent $V_A=V\cos\alpha$\quad $\alpha\downarrow$\quad $\cos\alpha\uparrow$

\noindent $V_A\uparrow$\quad $\therefore T>mg$

% ---- block 041-04 (原稿第41页) kind=method ----
\notefig[0.25]{figures/p041_e.png}
% 图内标注：直角三角形(斜边为杆/绳)，上端点 1 处标 $V_1$ 竖直向下及夹角 $\theta$，下端点 2 处标 $V_2$ 水平向右、夹角 $\theta$、虚线箭头
\[
V_1\cos\theta=V_2\sin\theta
\]

% ---- block 041-05 (原稿第41页) kind=note ----
\red{\struck{\unclear{难}}\ 沿绳方向速度相等}% 红笔；最左一字被红笔大叉划去(像“难”)；此红字是对上方各例的总结

% ---- block 041-06 (原稿第41页) kind=method ----
\notefig[0.25]{figures/p041_f.png}
% 图内标注：定滑轮(顶部横梁)，绳一端竖直挂物体 2(向上箭头)，另一端斜拉竖直轨道内的物体 1；物体 1 处标夹角 $\theta$，粗箭头向下标 $V_1$
\noindent $V_1\cos\theta=V_2$

\noindent $\theta\downarrow$\quad $V_1\cos\theta\uparrow$\quad $\therefore V_2\uparrow$\quad 2向上加速

\noindent $T>m_2g$

% ---- block 041-07 (原稿第41页) kind=method ----
\notefig[0.27]{figures/p041_g.png}
% 图内标注：斜杆(长 $L$，下端 $B$ 铰在斜面/地面左下，与水平夹角 $\theta$)，杆上端点 $A$；右下方块(高 $h$)以速度 $V$ 向右，其上角与杆接触，红箭头为垂直杆的分速度；红笔沿杆描出
\noindent \fcolorbox{red!75!black}{white}{\red{同轴同$\omega$}}% 红笔圈出，由 $A$ 处括号引出

\[
\left\{\begin{aligned}
&V\sin\theta=\omega\,\dfrac{h}{\text{\struck{\unclear{$\tan\theta$}}}\,\sin\theta}\\
&V_A=\omega L
\end{aligned}\right.
\]% 分母中最左一项被涂划(像 tanθ)，其后为 sinθ；原稿此处 ω 写得像 W

\noindent \red{同一转轴 $\omega$ 相同}

% ---- block 041-08 (原稿第41页) kind=method ----
\notefig[0.25]{figures/p041_h.png}
% 图内标注：一个五边形物体速度 $V_2$ 由绳(与水平夹角 $\alpha$)拉着；绳另一端为人，人速度 $V_1$ 水平向右(红箭头为沿绳/垂直绳分量)，虚线沿绳延长
\[
V_2=V_1\cos\alpha
\]

% ---- block 042-04 (原稿第42页) kind=method ----
\red{？}% 红笔问号，位于“两端圆周速度相等”左侧

\noindent 两端圆周速度相等\quad 一绳上 $T$、$x$、$v$、$a$、\unclear{$l$} 全相等% “a”后的一字像“1/I/l”

\noindent 绳不转动要平动

\notefig[0.25]{figures/p042_f.png}
% 图内标注：斜板(与水平夹角 $30^\circ$)，板上标 $V$(沿板向右上箭头)，板与竖直绳交点处标“板 $\frac12V$”(向上箭头)、$\frac{\sqrt3}{2}V$(水平向右箭头)；竖直绳下端方块，绳上标 $V_1$(向上箭头)，方块处水平向右箭头
\noindent $\dfrac{\sqrt3}{2}V$ 抵消\unclear{车}转动% 位于图下方右侧(第三字像“车”，疑为“轮”)

\begin{align*}
V+\tfrac12V&=V_1\\
V_1&=\tfrac32V\\
\therefore V_{\text{合}}&=\sqrt3\,V
\end{align*}


\section{关联速度·交点速度}
% ---- block 042-01 (原稿第42页) kind=derivation ----
\textbf{交点速度}% 页面左上角标题(后接虚线)

\notefig[0.3]{figures/p042_a.png}
% 图内标注：两条直线(其中一条水平线上标 $V_2\Delta t/\sin\alpha$)交于 $M$，$V_2$(小字)，$V_1\Delta t/\sin\alpha$、$V_1\Delta t$、向下箭头 $V_1$；$M$ 处向下箭头 $V_M$；虚线斜线末端箭头 $V_2$
\notefig[0.25]{figures/p042_b.png}
% 图内标注：三个矢量——向左箭头标 $V_2\Delta t/\sin\alpha$，向左下箭头标 $V_M$，向右下箭头标 $V_1\Delta t/\sin\alpha$
\[
V_M=\frac{\sqrt{\left(\dfrac{V_2\Delta t}{\sin\alpha}\right)^{2}+\left(\dfrac{V_1\Delta t}{\sin\alpha}\right)^{2}-2\left(\dfrac{V_1\Delta t}{\sin\alpha}\,\dfrac{V_2\Delta t}{\sin\alpha}\right)}}{\Delta t}
\]% CHECK: 余弦定理的交叉项按理应含 cos(夹角)，原稿括号内未见 cos 因子，照原样转录

% ---- block 042-02 (原稿第42页) kind=derivation ----
\notefig[0.25]{figures/p042_c.png}
% 图内标注：两个相互交叠的菱形(旋转45°的方块)，左边菱形上标 $V$(向左箭头)，右边菱形上标 $2V$(向右箭头)，交点标 $M$
小量分析，延长\ 找角。

\notefig[0.28]{figures/p042_d.png}
% 图内标注：两条45°交叉直线交于 $M$(沿一线标 $\sqrt2V\Delta t$)，左下角 $45^\circ$，底部虚线上标 $2V\Delta t$，右侧标 $V\Delta t$、$\frac{\sqrt2}{2}V\Delta t$、角 $\alpha$，由 $M$ 向下的箭头 $V_M$，数条虚线(延长线)
% 图上方的文字“小量分析，延长 找角。”在图右上方(已在上面转录)

\noindent 用余弦定理\quad 勾股定\unclear{理}% 该行位于页面右上，行末被照片边缘截断

\[
V_M=\frac{\sqrt{\left(\sqrt2V\Delta t\right)^{2}+\left(\dfrac{\sqrt2}{2}V\Delta t\right)^{2}}}{\Delta t}=\frac{\sqrt{10}}{2}V
\]% 式末“V”右下有一小圆圈，像句号“。”，也可能是下标0(V_0)，未转录


\section{关联速度·伸缩等速}
% ---- block 042-03 (原稿第42页) kind=method ----
\notefig[0.3]{figures/p042_e.png}
% 图内标注：左端小圆圈处标 $V_1$(向左箭头及虚线分量箭头)，细绳自左端圆圈斜向下到中间点、再斜向上连到右端小圆圈；右段绳上标 $V_3$(沿绳虚线箭头)；中间点处标 $V$(水平向右箭头)、$V_2$(斜向下箭头)
\begin{keybox}[伸缩等速]% 红笔框
\[
V_1+V_2=V_3\qquad \sum V_{\text{伸长}}=\sum V_{\text{缩短}}
\]
\end{keybox}


\section{关联速度·弹性绳}
% ---- block 042-05 (原稿第42页) kind=method ----
\textbf{弹性绳}\ —\ 弹簧/激光/光线/轻杆套环。

\notefig[0.28]{figures/p042_g.png}
% 图内标注：水平横线与竖直线、斜线(由竖线下端斜向上交于横线)；交点处黑字 $V_{\text{伸长}}$、$V_{\text{圆周}}$、$\omega$(弯箭头)；红字 $V\sin\theta$、$V$(向右箭头)、$V\cos\theta$，另有一小红字(像“沿”)

\noindent \underline{弹性物质无速度关联}。% 红笔下划线，并延伸为红箭头指向右侧红框

\begin{keybox}% 红笔框
\notefig[0.25]{figures/p042_h.png}
% 图内标注：竖直线与水平线、弹簧(锯齿线)斜连在右上方小圆环与左下方小球之间(弹簧长度与角度都在变)
长度角度都变化、无法分解。
\end{keybox}


\section{关联速度·双绳拉船}
% ---- block 042-06 (原稿第42页) kind=method ----
\textbf{双绳拉船。}

\notefig[0.2]{figures/p042_i.png}
% 图内标注：两只船 $A$、$B$(各带向外上方箭头)，用绳呈“Y”形系于下方物体 $C$，$C$ 上标 $V_C$(向上箭头)
\notefig[0.28]{figures/p042_j.png}
% 图内标注：圆上三点 $A$(左)、$B$(右)、$C$(下)，连 $AB$、$AC$、$BC$，$C$ 处红字角 $\alpha$，竖直线由 $C$ 向上到圆顶并标 $V_C$(箭头)
\noindent 连$AB$\ 过$\overset{C}{B}$作画圆\qquad $V_C=\dfrac{|AB|}{\sin\alpha}$% “B”上方另有一小字母“C”(旁注/插入)，“B”笔画有叠写

\[
AB^{2}=\sqrt{V_A^{2}+V_B^{2}-2V_AV_B\cos\alpha}
\]% CHECK: 原稿等号左边写作 AB^2(上标2)，右边带根号，应为 AB=\sqrt{…} 或 AB^2 无根号，照原样转录


\section{关联速度·接触关联}
% ---- block 042-07 (原稿第42页) kind=knowledge ----
\textbf{接触关联}（左侧花括号分三条）
\begin{itemize}
\item 点点接触\quad \underline{不分离 $\Leftrightarrow$ $V$大小方向相同}
\item 点面接触\quad \underline{不分离 $\Leftrightarrow$ 垂直于接触面}\unclear{…}% 行末被照片边缘截断，其后可能还有字；箭头写作“(=)”形(手写 ⇔)，按 $\Leftrightarrow$ 转录
\item 面面接触\quad \underline{不分离} $\Leftrightarrow$ \underline{方向速度相同}
\end{itemize}

% ---- block 042-08 (原稿第42页) kind=derivation ----
\notefig[0.25]{figures/p042_k.png}
% 图内标注：斜杆(下端小圆圈为转轴 $O$，与水平夹角 $\theta$，杆长 $L$)，杆上端 $A$，杆上沿杆箭头 $V_{B\parallel}$；杆与方块上角接触点 $B$，$V_B$(水平向右)、垂直杆的分速度 $V_{B\perp}$(斜向下箭头)、$M$(方块下沿处)、方块高 $h$
\begin{align*}
V_M&=V_B\\
V_{B\perp}&=V_B\sin\theta=\omega\cdot|OB|\\
\therefore\omega&=\frac{V\sin\theta}{|OB|}=\frac{V\sin\theta}{\dfrac{h}{\sin\theta}}=\frac{V\sin^{2}\theta}{h}\\
V_A&=L\cdot\omega=\frac{LV\sin^{2}\theta}{h}
\end{align*}% 原稿 ω 写得像 W


\section{关联速度·滑轮组位移关联}
% ---- block 042-09 (原稿第42页) kind=problem ----
\red{？}% 红笔大问号，位于该题图左上方

\notefig[0.3]{figures/p042_l.png}
% 图内标注：竖直墙面(左)、水平细线与两个小滑轮(圆圈)，水平段向左箭头；斜面(倾角 $\alpha$)上有物块 $B$(标 $m$，沿斜面向下的箭头)，红笔：$B$ 红字、红线连接滑轮与物块、红字 $X_B$ 及红色三角；斜面下方 $m$、$A$、$\alpha$，虚线；图下方“光滑”
\noindent $A$ 位移为 $x$ 时，求 $S_B$、$V_A$

\noindent 对$B$ 先向左，再沿斜面向下。

\notefig[0.2]{figures/p042_m.png}
% 图内标注：位移矢量图——水平向左的 $x$(虚线)、沿斜面向下的 $x$、两段夹角红字 $\alpha/2$ 等，红粗线为合位移(斜向下箭头)
\[
S_B=2x\sin\frac{\alpha}{2}
\]

\[
mgx\sin\alpha=\frac12mV_x^{2}+\frac12m\left(V_x\cdot2\sin\frac{\alpha}{2}\right)^{2}
\]

\[
\therefore V_x=\sqrt{\frac{2gx\sin\alpha}{3-2\cos\alpha}}
\]


\section{关联速度·几何关联}
% ---- block 043-01 (原稿第43页) kind=knowledge ----
\begin{keybox}[几何关联]% 红笔框(页面左上角)，框内两行
$V$ 是三角形直角边
\end{keybox}

% ---- block 043-02 (原稿第43页) kind=note ----
\red{受力要整体法}% 红色大字，位于页面右上角

% ---- block 043-03 (原稿第43页) kind=problem ----
\begin{problem}[{[1]}]
\notefig[0.3]{figures/p043_a.png}
% 图内标注：图左上角红框“几何关联 / V是三角形直角边”(已在上一块转录)；楔形斜面(左下角 $\theta$)，斜面上顶着一根竖直杆，杆顶标 $M$(重物)；楔形右侧标 $m$，水平向左箭头 $F$；红字 $V_M$(竖直向下，右侧)、$V_m$(水平，楔块下方红线箭头)
① 为将重物顶起，求 $F_{\min}$

\noindent $F_N=\dfrac{Mg}{\cos\theta}$\qquad $F_N\sin\theta=\dfrac{Mg}{\cos\theta}\sin\theta$

\noindent 将小轮推高地面 $h$ 处，突然撤去 $F$，求当小轮到水平面时斜面速度

\noindent ②
\[
\left\{\begin{aligned}
&V_M=V_m\tan\theta\\
&Mgh=\tfrac12MV_M^{2}+\tfrac12mV_m^{2}
\end{aligned}\right.
\]
\end{problem}

% ---- block 043-04 (原稿第43页) kind=problem ----
\begin{problem}[{[2]}]
\notefig[0.3]{figures/p043_b.png}
% 图内标注：竖直墙(左，顶端标 1)与斜面体(倾角 $30^\circ$，斜面上方标 2、$u=0$)之间夹一个小球(半径 $R$，标 $m$、两处 $R$)，球心到切点的虚线；斜面体右侧竖直边标 $M$，向左箭头 $F$，右侧红色小箭头 $V_1$，底边红箭头 $V_2$，右下标 $u=0$
撤去 $F$ 后小球和斜面体做匀加速直线运动，直到小球落地。

\noindent ① 求 $F$\quad ② 球落地时求 $V_1$、$V_2$

\noindent $F=N=\dfrac{\sqrt3}{3}mg$

\noindent $mgh=\tfrac12mV_1^{2}+\tfrac12mV_2^{2}$

\noindent $h_0=\sqrt3R$\quad $h=h_0-R=(\sqrt3-1)R$

\noindent $x_1=x_2\tan\theta$

\noindent $V_2=\sqrt3V_1$

\noindent $V_1=\sqrt{\dfrac{(\sqrt3-1)gR}{2}}$\quad $V_2=\sqrt{\dfrac{3(\sqrt3-1)gR}{2}}$
\end{problem}

% ---- block 043-05 (原稿第43页) kind=problem ----
\begin{problem}[{[3]}]
\notefig[0.3]{figures/p043_c.png}
% 图内标注：左侧楔形物块 $A$(顶角处标 $m$，下角 $37^\circ$、$\alpha$，红字 $V_B$ 水平向右箭头在顶端，红字 $V_A$ 竖直向下箭头在左侧)，紧靠的方块 $B$(标 $2m$)，弹簧(标 $k$)连接 $B$ 与右侧竖直墙；图下标 $\mu=0$；图左上角有题号 [3]
$B$ 将弹簧压缩 $x_0$ 后 $A$ 上滑，弹簧恢复原长时，$B$ 速度为 $V_B$。

\noindent ① 求刚恢复原长时 $V_A$。

\noindent $V_B=V_A\tan\alpha$\qquad $V_A=\dfrac43V_B$

\notefig[0.25]{figures/p043_d.png}
% 图内标注：矢量图——向上箭头 $V_A$，向左箭头 $V_B$，夹角 $37^\circ$，虚线(斜向左上)
\noindent ② 求弹簧压缩量为 $x_0$ 时具有的弹性势能

\noindent 弹簧原长时，$h_A=\dfrac{x_0}{\tan\alpha}=\dfrac43x_0$

\noindent $E_p=\tfrac12mV_A^{2}+\dfrac{2}{2}mV_B^{2}+mgh=\dfrac{17}{9}mV_B^{2}+\dfrac43mgx_0$

\noindent ③ 求两物体动能最大时，弹簧压缩量

\noindent 设$AB$相互作用力$F$
\[
\left\{\begin{aligned}
&a_A=a_B=0\\
&F\sin\alpha=mg\\
&F\cos\alpha=kx
\end{aligned}\right.\qquad x=\frac{mg}{k\tan\alpha}=\frac{4mg}{3k}
\]
\end{problem}


\section{平抛运动·基本公式}
% ---- block 044-01 (原稿第44页) kind=formula ----
\notefig[0.3]{figures/p044_a.png}
% 图内标注：斜面(下端标 $\theta$、$\alpha$)，抛体的抛物线轨迹；上端抛出点处 $V_0$(向左箭头)、$V\sin\theta$(向上箭头)，沿斜面 $V\cos\theta$；加速度分解 $g\sin\theta$、$g$(向下箭头)、$g\cos\theta$(斜向下箭头)，轨迹中部有小球及 $\theta$ 标记
\begin{align*}
x&=V_0t & V_x&=V_0\\
y&=\tfrac12gt^2 & V_y&=gt\\
S&=\sqrt{x^2+y^2} & V&=\sqrt{V_x^2+V_y^2}\\
\tan\alpha&=\frac{y}{x}=\frac{gt}{2V_0} & \tan\theta&=\frac{V_y}{V_x}=\frac{gt}{V_0}
\end{align*}

% ---- block 044-02 (原稿第44页) kind=note ----
\notefig[0.25]{figures/p044_b.png}
% 图内标注：“类上抛”“匀加速”两行字，左侧一个向左上箭头，另有一条斜向左下箭头由“匀加速”指向左下(指向左侧斜面图)
\notefig[0.25]{figures/p044_c.png}
% 图内标注：L 形箭头——水平向右箭头旁写“匀速直线”，竖直向下箭头下写“自由落体”
\noindent $\nwarrow$ 类上抛\quad 匀加速$\swarrow$% 二者在页面右上方，箭头指向左侧斜面图

\noindent $\rightarrow$ 匀速直线\quad $\downarrow$ 自由落体

\noindent \red{抛体\unclear{重}分解}

\noindent \red{平抛\unclear{有}两个直角三角形\ 一个反向延长线}

\noindent \red{斜抛\unclear{有}从顶部的两个平抛}% 以上三行均为红笔；“有”字辨认不清(像“考”)


\section{平抛运动·斜面平抛}
% ---- block 044-03 (原稿第44页) kind=formula ----
\textbf{① 落点时间}
\[
\tan\theta=\frac{gt}{2V_0}=\frac12\tan\alpha\qquad
\tan\theta=\frac{\frac12gt^2}{V_0t}\qquad
t=\frac{2V_0\tan\theta}{g}
\]

% ---- block 044-05 (原稿第44页) kind=formula ----
\textbf{③ 离斜面最远}
\[
h=\frac{(V_0\sin\theta)^2}{2g\cos\theta}\qquad
t=\frac{V_0\tan\theta}{g}\qquad
V_{\parallel}=V_0\cos\theta+g\sin\theta\,t
\]
\noindent 半程$\nearrow$% 小字“半程”带箭头，指向 t 式

% ---- block 044-06 (原稿第44页) kind=method ----
\textbf{④}

\notefig[0.25]{figures/p044_d.png}
% 图内标注：斜面上方的抛物线轨迹，沿斜面方向分成两段，分别标 1、2，右端标 $S_1$(带向左箭头)
\noindent $\dfrac{S_1}{S_2}=\dfrac13$\qquad $V_{y_0}=0$

\notefig[0.25]{figures/p044_e.png}
% 图内标注：$V$-$t$ 图像(纵轴 $V$，横轴 $t$)，过原点的直线下三角形按相等时间段分成面积 1、3、5 三块
\noindent $V_{y_0}\neq0$ 时\quad 若为

\notefig[0.25]{figures/p044_f.png}
% 图内标注：$V$-$t$ 图像(上部为过原点直线下的三角形，分成 1、3 两块；下部为两个相等矩形，各标 $S_0$)
\noindent \underline{真分数不等式。}% 红线之外为黑笔下划线

\[
\frac{1+S_0}{3+S_0}>\frac13
\]
\[
\therefore 1>\frac{S_1}{S_2}>\frac13
\]% 末式下方有波浪线

% ---- block 045-01 (原稿第45页) kind=knowledge ----
\notefig[0.25]{figures/p045_a.png}
% 图内标注：斜面，抛出点处水平箭头 $V_0$，几条抛物线落到斜面上，标 $2\alpha$、$\alpha$ 等角
\noindent 位偏 = 从斜到斜 = 位偏角为倾斜角\qquad $\tan\alpha=\dfrac{gt}{2V_0}$

% ---- block 045-02 (原稿第45页) kind=method ----
\notefig[0.3]{figures/p045_b.png}
% 图内标注：同一抛出点，两个水平初速度 $V_0$、$2V_0$ 的抛物线落到斜面上，红色直线连抛出点与两个落点，落点处标 $\alpha_1$、$\alpha_2$，右下标 $S_1$ 并写“求 $S_1/S_2$”
\noindent 求 $\dfrac{S_1}{S_2}$

\noindent 讨论\quad ① 都落在斜面上

\noindent $\tan\alpha=\dfrac{gt}{2V_0}$\qquad $t=\dfrac{2V_0\tan\alpha}{g}\propto V_0$\qquad \fbox{$t_1:t_2=1:2$}

\noindent $S=V_0t\propto V_0^2$\qquad \fbox{$\dfrac{S_1}{S_2}$ 则为 $1:4$}% 方框内“则”字辨认略不确定

\noindent ② 都落在平面

\noindent $h=\tfrac12gt^2$\quad \fbox{$t_1=t_2$}\quad $S=V_0t\propto V_0$\quad \fbox{$\dfrac{S_1}{S_2}=1:2$}

\noindent ③ 一个斜面一个平面

\noindent $\alpha_1>\alpha_2$

\noindent $\therefore\tan\alpha_1>\tan\alpha_2$

\noindent $\therefore\dfrac{gt_1}{2V_0}>\dfrac{gt_2}{2V_0}$

\noindent \fbox{$\begin{array}{l}\therefore\dfrac12<\dfrac{t_1}{t_2}<1\\[4pt]\therefore\dfrac14<\dfrac{S_1}{S_2}<\dfrac12\end{array}$}% 该方框位于 ③ 推导式右侧

% ---- block 045-03 (原稿第45页) kind=method ----
\notefig[0.3]{figures/p045_c.png}
% 图内标注：两个对称斜面构成 V 形，各高 $h$、水平长 $2h$，左斜面顶端抛出 $V$(水平箭头)，抛物线落到谷底 $O$ 点；$O$ 处标 $\theta_1$、$\theta_2$，向下箭头
\noindent $t_{\text{最短}}$\quad $t=\sqrt{\dfrac{2h}{g}}$

\noindent 找界点：在 $O$ 点\quad $\theta_1<45^\circ$\qquad $\tan\theta_2=2\tan\theta_1=\dfrac{2y}{x}=1$

\noindent $\theta_1+\theta_2<90^\circ$\qquad 故无法垂直打在斜面 (若末速度反向\unclear{延}长线过位移中点。)% 括号内文字位于页面右侧，第一行末“过”字靠近页边，第二行为“位移中点。)”

% ---- block 045-04 (原稿第45页) kind=problem ----
\notefig[0.3]{figures/p045_d.png}
% 图内标注：左上小圆圈出发水平箭头 $V_0$(抛出点)；斜面(高 $h$，带花括号)，斜面上沿斜面向下箭头 $V_1$；斜面底端连水平面，水平面上标 $x$、$\mu$；图下斜写小字“如履平地”(第二字像“覆”)
\noindent 如\unclear{履}平地% 斜写于图下方

\noindent 求能走多远
\[
\left\{\begin{aligned}
V_1&=V_0\cos\theta+gt\sin\theta\\
\tan\theta&=\frac{gt}{2V_0}\\
h&=H-\tfrac12gt^2\\
\tfrac12mV_1^2&=mgh-\mu mg\cos\theta\,\frac{h}{\sin\theta}-\mu mgx
\end{aligned}\right.
\]


\section{平抛运动·重力功率}
% ---- block 044-04 (原稿第44页) kind=formula ----
\textbf{② 重力功率}
\[
\left\{\begin{aligned}
\bar P_G&=\frac{mgy}{t}=\frac{mg\cdot\frac12gt^2}{t}=\frac12mg^2t\\
P_G&=mgV\cos\alpha=mgV_y=mg^2t
\end{aligned}\right.
\qquad
\left\{\begin{aligned}
\bar V&=V_{\frac t2}=\frac{V_0+V}{2}\\
\bar P&=P_{\frac t2}=\frac{P_0+P}{2}
\end{aligned}\right.
\]% 一条细弧线自 mgVcosα 附近向左下延伸至 ③ 的 h 式附近(指示线)


\section{平抛运动·轨迹方程}
% ---- block 044-07 (原稿第44页) kind=formula ----
\textbf{⑤ 轨迹方程}\quad $y=\dfrac{g}{2V_0^2}x^2$\quad (用于联立，建系) 找交点。
\[
\left\{\begin{aligned}x&=V_0t\\y&=\tfrac12gt^2\end{aligned}\right.
\]


\section{平抛运动·偏角}
% ---- block 044-08 (原稿第44页) kind=formula ----
\textbf{⑥ 偏角}
\[
\left\{\begin{aligned}
\tan\theta&=\frac{gt}{V_0}\quad\text{速偏}\\
\tan\alpha&=\frac{gt}{2V_0}\quad\text{位偏}
\end{aligned}\right.
\]

\notefig[0.25]{figures/p044_g.png}
% 图内标注：水平线上两个反向箭头，左边标 $V_1$、右边标 $V_2$，中间竖直虚线；箭头旁标 1(向左)、2(向右)
\noindent 当 1、2 速度垂直时求 $t$
\[
\tan\theta_1\cdot\tan\theta_2=1
\]
\[
\tan\theta_1=\frac{gt}{V_1}\qquad\tan\theta_2=\frac{gt}{V_2}\qquad t=\frac{\sqrt{V_1V_2}}{g}
\]


\section{平抛运动·抛上另一斜面}
% ---- block 044-09 (原稿第44页) kind=method ----
\textbf{⑦ 抛上另一个斜面}

\notefig[0.25]{figures/p044_h.png}
% 图内标注：$V_0$(水平向右箭头)，抛物线落到斜面上，斜面倾角 $\theta$，落点处速度箭头及夹角 $\varphi$
\noindent 若垂直，$\tan\varphi=\dfrac{gt}{V_0}=\dfrac{1}{\tan\theta}$

\[
\bar P=\frac{mgy}{t}=\frac{mg\cdot\frac12gt^2}{t}=\frac{mgV_0\cot\theta}{2}=mgV_y
\]% CHECK: 末项 mgV_y 与前一项 mgV_0cotθ/2 相差系数 1/2，疑漏写或 V_y 指平均值，照原样转录；“P̄=”之后有一小块被涂白

\notefig[0.25]{figures/p044_i.png}
% 图内标注：$V_0$(水平向右箭头)，斜面倾角 $\theta$，轨迹与斜面“最短距离”处(标“最短距离”并带引线)，速度与斜面夹角 $90^\circ-\theta$
\[
\frac{gt}{2V_0}=\tan(90^\circ-\theta)=\frac{1}{\tan\theta}
\]

\notefig[0.25]{figures/p044_k.png}
% 图内标注：左侧花括号标 $H$(总高度)；$V_0$ 水平向右箭头，抛物线落到倾角 $\theta$ 的斜面上，落点速度方向 $\theta$，虚线辅助线；两条细长弧线由图中两段高度分别指向右侧式中 $V_0t\tan\theta$ 与 $\frac12gt^2$ 的下划线
\[
\left\{\begin{aligned}
H&=\underline{V_0t\tan\theta}+\underline{\tfrac12gt^2}\\
\frac{1}{\tan\theta}&=\frac{gt}{V_0}
\end{aligned}\right.
\qquad V_0=\sqrt{gH}
\]% CHECK: 由左式联立得 H=(V_0^2/g)(1+cot^2θ/2)，与 V_0=√(gH) 不一致，照原样转录


\section{平抛运动·偏角之差(不单调)}
% ---- block 044-10 (原稿第44页) kind=method ----
\textbf{⑧}

\notefig[0.25]{figures/p044_j.png}
% 图内标注：斜面(倾角 $\varphi$)，上方箭头轨迹，标 $\alpha$、$\theta$、$\varphi$ 等角，水平虚线
\[
\tan\theta=\tan(\alpha-\varphi)=\frac{\tan\alpha-\tan\varphi}{1+\tan\alpha\tan\varphi}
\]
\[
\tan\alpha=2\tan\varphi
\]
\[
\therefore\tan\theta=\frac{1}{2\tan\varphi+\dfrac{1}{\tan\varphi}}
\]

\noindent 不单调
\notefig[0.25]{figures/p044_l.png}
% 图内标注：坐标轴(纵轴带箭头)与一条过原点的波形曲线：原点左侧先下凹，右侧上凸后下降，旁写“不单调”


\section{平抛运动·无碰撞}
% ---- block 044-11 (原稿第44页) kind=method ----
\noindent 无碰撞 = 速偏 $\Leftrightarrow$ $V\parallel$接触面% “接触面”三字写得较连，辨认为“接触面”

\notefig[0.25]{figures/p044_m.png}
% 图内标注：$\Rightarrow$ 水平箭头，虚线轨迹，到斜面上的小物体(管口)，弯箭头；图下标“平行”
\notefig[0.25]{figures/p044_n.png}
% 图内标注：“0$\Rightarrow$”水平箭头，虚线轨迹，进入半圆形碗(右上是碗口)，向下箭头及虚线；图下标“切向”
\noindent 平行\qquad 切向

\notefig[0.25]{figures/p044_o.png}
% 图内标注：$\Rightarrow V_0$ 水平箭头，虚线轨迹落到半圆弧上，圆弧上标 $60^\circ$、$30^\circ$、$120^\circ$(像“160°”)，末端速度箭头 $V$
\noindent 刚好相切无碰

\[
\left\{\begin{aligned}
\tan30^\circ&=\frac{gt}{V_0}\\
\frac{\sqrt3}{2}R&=V_0t
\end{aligned}\right.
\]% CHECK: “R”前分数的分子写得像“3”(根号不明)，按 √3/2 转录

\noindent \struck{$\therefore V_0=3\sqrt5\unit{m/s}$}% 整行被横线划去

\noindent $\Rightarrow V_0$


\section{圆周运动·竖直面临界}
% ---- block 045-05 (原稿第45页) kind=summary ----
\notefig[0.25]{figures/p045_e.png}
% 图内标注：每行左侧的示意小图——第一行红笔：竖线上端带小圆(绳)、半圆弧∪(槽)、小方框(筒)；第二行：红色竖线带小圈(杆)，“环”“管”二字各被红圈圈出；第三行：黑色拱形∩
\noindent $mg+T=\dfrac{mV^2}{r}$% 位于表格上方、第一行方程栏

\noindent \red{$V<\sqrt{gr}$ 或 $\omega<\sqrt{\dfrac{g}{r}}$ 则对筒槽绳压力为0 即为最小值}% 红笔，位于第一行上方

\begin{center}
\small
\begin{tabular}{@{}l l l c c c c@{}}
\toprule
类型 & 受力特点 & 最高点方程 & 最高 & 最低 & 上下两点压力差$\Delta N$ & 向心力上下差\\
\midrule
绳\ 槽\ 筒 & 无法提供支持力、给压力 & $mg+T=\dfrac{mV^2}{r}$ & $V=\sqrt{gR}$ & $V=\sqrt{5gR}$ & $6mg$ & $4mg$\\[4pt]
杆\ 环\ 管 & 双向力 & $mg\pm F_N=\dfrac{mV^2}{r}$ & $0$ & $V=\sqrt{4gR}$ & $6mg$ & $4mg$\\[4pt]
拱 & 支持力 & $mg-N=\dfrac{mV^2}{r}$ & $0$ & 无 & 无 & 无\\
\bottomrule
\end{tabular}
\end{center}% 原稿表格左侧有一个大花括号括住三行；“类型/受力特点/最高点方程”三个表头为转录者补写的栏名(原稿无)；第一行的方程 mg+T=mV^2/r 在原稿中位于该行上方；第三行方程前有一小块被涂白

% ---- block 045-06 (原稿第45页) kind=method ----
\noindent ☆ 不脱离、\underline{绷直}

\notefig[0.2]{figures/p045_f.png}
% 图内标注：圆(半径 $R$)，圆心到右侧等高点的水平虚线、到最低点的竖直虚线；圆上标 $V_2$(最高点)、$V_1$(与圆心等高的右侧点)；最低点有一个小圆圈(出发点)
\noindent $V\le V_1$ 或 $V\ge V_2$

\noindent $V\le\sqrt{2gR}$\qquad $V\ge\sqrt{5gR}$

% ---- block 046-03 (原稿第46页) kind=method ----
\textbf{技巧}

\notefig[0.3]{figures/p046_a.png}

\[
\begin{cases}
mg\,2R+W_{\text{其他}}=\dfrac{1}{2}mV_{\text{下}}^{2}-\dfrac{1}{2}mV_{\text{上}}^{2}\\[4pt]
N_{\text{上}}+mg=\dfrac{mV_{\text{上}}^{2}}{R}\\[4pt]
N_{\text{下}}-mg=\dfrac{mV_{\text{下}}^{2}}{R}
\end{cases}
\]
\[ \therefore\ N_{\text{下}}-N_{\text{上}}=6mg+\frac{2W_{\text{其他}}}{R} \]
\[ \Delta N=6mg+\frac{2W_{\text{其他}}}{R} \]
\[ \text{实际}\,\Delta N=\text{理想}\,\Delta N+2\,\frac{W_{\text{其他}}}{R} \]
\[ \text{偏差}\,\Delta N\,\frac{R}{2}=W_{\text{其他}}\qquad W_{\text{其他}}=\frac{\Delta N_{\text{偏差}}}{2}R\qquad \Delta N_{\text{偏差}} \]


\section{圆周运动·圆心参考系}
% ---- block 046-01 (原稿第46页) kind=knowledge ----
$V$为相对圆心的速度，而非单单相对地面

圆心参考系．\ $F_n=\sum F_{\text{指向圆心}}-\sum F_{\text{背离圆心}}$\ （\unclear{分解}）

\[ F_n=\frac{mv^2}{R}=\frac{2E_k}{R} \]

% ---- block 089-03 (原稿第89页) kind=problem ----
\noindent 圆周运动以圆心为参考系\qquad $V$ 为相对于圆心的速度。
\begin{problem}
\notefig[0.25]{figures/p089_b.png}
当B\underline{第一次回到$A$}正下方时，求细线拉力大小
% 下划线画在“第一次回到A”下面
\begin{solution}
完全弹性碰撞\qquad $\Delta V_{\text{前}}=\Delta V_{\text{后}}$\qquad $T-mg=\dfrac{m\Delta V^2}{L}=\dfrac{mV_0^2}{L}$
% “ΔV后”的 Δ 手写形似 0
\end{solution}
\end{problem}


\section{圆周运动·传动}
% ---- block 046-02 (原稿第46页) kind=knowledge ----
传动\ $\begin{cases}\text{同点同}\,V\\ \text{同轴同}\,\omega\end{cases}$


\section{运动合成·相对速度}
% ---- block 046-04 (原稿第46页) kind=knowledge ----
\textbf{运动合成}\ \ {\footnotesize $P$，$I$，$V$，$a$，$X$}

\[ ab^{\text{对}}=a_{\text{对地}}-b_{\text{对地}} \]
% CHECK: 公式左端"αb"及其右上角的"对"辨认不清；右端"对地"的"对"写在"地"字上方；"运动合成"右上方有小字 P、I、V、a、X（疑为动量、冲量、速度、加速度、位移皆为相对量）

\notefig[0.5]{figures/p046_b.png}
% 图 p046_b 内含(黑字)公式 V雨车=V雨地−V车地 及矢量三角形、各标注

\[ V_{\text{雨车}}=V_{\text{雨地}}-V_{\text{车地}} \]

\notefig[0.3]{figures/p046_d.png}

\[ V_{\text{刀玻}}=V_{\text{刀地}}-V_{\text{玻地}} \]


\section{小船过河}
% ---- block 046-05 (原稿第46页) kind=method ----
\textbf{小船过河}\quad ① 最快\ $t=\dfrac{d}{V_{\text{船}}}$\quad ② 最近\ $V_{\text{船}}>V_{\text{水}}$\ 垂直过河

$V_{\text{船}}<V_{\text{水}}$\quad $V_{\text{水}}$为斜边，$V_{\text{船}}$画圆找切线\quad $\sin\alpha=\dfrac{d}{L}=\dfrac{V_{\text{船}}}{V_{\text{水}}}$

\notefig[0.45]{figures/p046_e.png}

\notefig[0.3]{figures/p046_f.png}

甲乙$V$相等\unclear{向}P，在NP上\unclear{相}遇，渡河时间一致

\notefig[0.25]{figures/p046_g.png}

$V_2$相对甲

\notefig[0.35]{figures/p046_h.png}

船头去\unclear{时}垂直河岸，归\unclear{时}路线垂直河岸，

\[ t_1=\frac{d}{V_{\text{船}}}\qquad t_2=\frac{d}{\sqrt{V_{\text{船}}^{2}-V_{\text{河}}^{2}}}\qquad \frac{t_1}{t_2}=k\qquad k=\sqrt{1-\left(\frac{V_{\text{水}}}{V_{\text{船}}}\right)^{2}} \]
% CHECK: t_2 分母根号内 V船² 与 V河² 之间的符号原稿像"="（疑为"−"的笔误），此处按"−"录入；t_2 中写 V河、k 中写 V水，照原稿


\section{曲线运动·条件与速度变化}
% ---- block 068-01 (原稿第68页) kind=knowledge ----
% 原稿：页面最上方一行（位于“Date. NO.”表头上方），以及表头下方右侧两行、其下两处批注
$a$\ $V$锐角\ $V\uparrow$\qquad $a$\ $V$钝角\ $V\downarrow$\qquad $a$\ $V$\unclear{垂直}\ 速度大小不变，方向改变 % CHECK: 顶行三项中 a 与 V 之间似无符号（应为夹角关系），第三项“垂直”二字不确定

$a\perp V$\quad 改变$V$方向\quad 法向加速度

$a\parallel V$\quad 改变$V$大小\quad 切向加速度

$a$、$V$不同方向

力与速度夹轨迹\quad 凹侧方向偏向力


\section{章节提纲}
% ---- block 068-02 (原稿第68页) kind=summary ----
% 版面：左侧一列为各部分标题（曲线运动；运动的合成与分解，其下一行小字“合成、分解”；抛体运动，其下“分解”；圆周运动，其下一行小字），右侧为对应条目
\begin{itemize}
  \item \textbf{曲线运动}
  \begin{itemize}
    \item 曲线运动、骑马射箭、转台投篮、风中骑车、\unclear{滑}板运动
  \end{itemize}
  \item \textbf{运动的合成与分解}\\ 合成、分解
  \begin{itemize}
    \item 运动的合成与分解，\quad 等时性\quad 独立性\quad 等效性
    \item 小船渡河模型、
    \item 关联速度模型、 % 图及批注见 068-03
  \end{itemize}
  \item \textbf{抛体运动}\\ 分解
  \begin{itemize}
    \item （类）平抛运动
    \item （类）斜抛运动
    \item 落点有约束的平抛运动
    \item 平抛运动的临界问题
    \item 斜抛运动问题。
  \end{itemize}
  \item \textbf{圆周运动}\\ \unclear{受力}
  \begin{itemize}
    \item 线速度\quad \underline{角速度}\quad \underline{周期}、\underline{转速}\quad 向心加速度\quad 向心力
    \item 离心运动（$F_n<m\omega^2 r$）、和向心运动（$F_n>m\omega^2 r$） % 原稿两式分别写在“离心运动”“和向心运动”上方
    \item $\underbrace{\text{转杆}}_{\text{同}\,\omega}$、$\underbrace{\text{皮带传动、齿轮传动、摩擦传动}}_{\text{同}\,v}$、$\underset{\text{同}\,\omega}{\text{同轴传动}}$ % 原稿批注写在各项下方：转杆下波浪线+“同ω”；皮带/齿轮/摩擦传动下长波浪线+“同v”；同轴传动下“同ω”
  \end{itemize}
\end{itemize}


\section{关联速度模型}
% ---- block 068-03 (原稿第68页) kind=diagram ----
\textbf{关联速度模型}

\notefig[0.28]{figures/p068_a.png}

\notefig[0.28]{figures/p068_b.png}

沿绳垂绳，沿杆垂杆

力和速度的分解不同


\section{水平面圆周运动·飞机转弯}
% ---- block 069-01 (原稿第69页) kind=method ----
\textbf{水平面内的匀速圆周运动。}

\textbf{飞机水平转弯}

\notefig[0.25]{figures/p069_a.png}

\[ F_n=mg\tan\theta \]


\section{水平面圆周运动·火车转弯}
% ---- block 069-02 (原稿第69页) kind=method ----
\textbf{火车转弯}

\notefig[0.25]{figures/p069_b.png}

\[ F_n=mg\tan\theta \]


\section{水平面圆周运动·飞车走壁}
% ---- block 069-03 (原稿第69页) kind=method ----
\textbf{飞车走壁}

\notefig[0.4]{figures/p069_c.png}

\[ F_n=mg\tan\theta \]


\section{水平面圆周运动·汽车转弯}
% ---- block 069-04 (原稿第69页) kind=method ----
\textbf{汽车转弯}

\notefig[0.25]{figures/p069_d.png}

\[ f=F_n \]


\section{水平面圆周运动·圆筒内壁}
% ---- block 069-05 (原稿第69页) kind=method ----
\notefig[0.25]{figures/p069_e.png}

\begin{align*}
& mg=f\le\mu F_N\\
& F_N=m\omega^2 r\\
& \omega\ge\sqrt{\dfrac{g}{\mu r}}
\end{align*}
转速越大越不易掉下


\section{水平面圆周运动·水平转台}
% ---- block 069-06 (原稿第69页) kind=method ----
\textbf{水平转台}

\notefig[0.25]{figures/p069_f.png}

\[ F_n=m_Bg \]


\section{水平面圆周运动·圆锥摆}
% ---- block 069-07 (原稿第69页) kind=method ----
\textbf{圆锥摆}

\notefig[0.3]{figures/p069_g.png}

\begin{align*}
& T\cos\theta=mg\\
& T\sin\theta=mg\tan\theta=F_{\text{向}}
\end{align*}

\textbf{单个圆锥摆}
\begin{align*}
F_n&=mg\tan\theta & v&=\sqrt{ar}=\sqrt{g\tan\theta\sin\theta}\quad\text{与}\theta\text{正相关}\\ % CHECK: 根号内疑似漏写 L（r=L sin θ 时 v=sqrt(gL tanθ sinθ)），照原样转录
r&=L\sin\theta & \omega&=\sqrt{\dfrac{a}{r}}=\sqrt{\dfrac{g}{H}}\\
T&=mL\omega^2=\dfrac{mg}{\cos\theta} & a_{\text{向}}&=g\tan\theta=\omega^2r=\dfrac{v^2}{r}
\end{align*}

% ---- block 069-08 (原稿第69页) kind=derivation ----
% 左侧页边小标题为“水平转盘、单列”（第二行两字不确定），其右为下列推导
\textbf{水平转盘、\unclear{单列}}
\[ mg\tan\theta=m\omega^2L\sin\theta=\dfrac{m\cdot4\pi^2}{T^2}L\sin\theta,\qquad \omega=\sqrt{\dfrac{g}{L\cos\theta}}=\sqrt{\dfrac{g}{H}} \]
\[ T=2\pi\sqrt{\dfrac{L\cos\theta}{g}}=2\pi\sqrt{\dfrac{H}{g}} \]


\section{水平面圆周运动·双圆锥摆}
% ---- block 069-09 (原稿第69页) kind=method ----
\textbf{双圆锥摆}

\notefig[0.3]{figures/p069_h.png}

$\omega$相等，$V$和$a$、$\propto r$，

竖直方向合力为0


\section{竖直面圆周运动·轻绳模型}
% ---- block 070-01 (原稿第70页) kind=method ----
\textbf{竖直圆周运动}

\textbf{轻绳模型}

\notefig[0.3]{figures/p070_a.png}

\begin{align*}
& mg+F_T=\dfrac{mv^2}{r}\\
& F_T=0\qquad mg=\dfrac{mv^2}{r}\\
& v=\sqrt{gr}\\
& F=\dfrac{mv^2}{r}-mg
\end{align*}

无支持力

绳力只能向下。

\[ V_{\text{过上}}=\sqrt{gr} \]


\section{竖直面圆周运动·轻杆模型}
% ---- block 070-02 (原稿第70页) kind=method ----
\textbf{轻杆模型}

\notefig[0.3]{figures/p070_b.png}

弹力可向下\quad 可向上\quad 也为0

\unclear{当}\ $V=0$\quad $F_n=0$\quad $F_N=mg$

$V_{\text{过上}}=0$

\notefig[0.3]{figures/p070_c.png}

\begin{align*}
& mg\pm F_N=\dfrac{mv^2}{r}\\
& F=\dfrac{mv^2}{r}-mg\\
& mg-(-F)=\dfrac{mv^2}{r}
\end{align*}


\section{水平面圆周·圆锥内壁}
% ---- block 076-01 (原稿第76页) kind=method ----
\notefig[0.25]{figures/p076_a.png}
\[
\begin{aligned}
&N\sin\theta = mg\\
&N\cos\theta = mg\cot\theta\\
&a_{\text{向}} = g\cot\theta = \omega^2 r = \frac{v^2}{r}
\end{aligned}
\qquad\qquad
\begin{aligned}
&v=\sqrt{ar}=\sqrt{g\cot\theta\,r}\\
&\omega=\sqrt{\frac{a}{r}}=\sqrt{\frac{g}{H}}\\
&T=2\pi\sqrt{\frac{H}{g}}
\end{aligned}
\]
% CHECK: 图中 H 的含义不明（标在水平箭头旁），ω=√(g/H)、T=2π√(H/g) 与 a=g cotθ 的关系需核对


\section{水平面圆周·半球碗内壁}
% ---- block 076-02 (原稿第76页) kind=method ----
\notefig[0.3]{figures/p076_b.png}
\[
\begin{aligned}
&N\cos\theta = mg\\
&N\sin\theta = mg\tan\theta = F_n\\
&r = R\sin\theta
\end{aligned}
\qquad\qquad
\begin{aligned}
&v=\sqrt{ar}=\sqrt{g\tan\theta\sin\theta}\quad\text{正相关}\\
&\omega=\sqrt{\frac{a}{r}}=\sqrt{\frac{g}{H}}=\sqrt{\frac{g\tan\theta}{L\sin\theta}}\\
&T=2\pi\sqrt{\frac{H}{g}}
\end{aligned}
\]
% CHECK: v=√(g tanθ sinθ) 手写如此，按 r=R sinθ 似漏写 R；ω 末式分母手写像 L，按 r=R sinθ 应为 R

% ---- block 076-03 (原稿第76页) kind=derivation ----
\notefig[0.3]{figures/p076_c.png}
\textbf{若光滑}
\begin{align*}
&N\cos\theta = mg\\
&N\sin\theta = mg\tan\theta = ma\\
&H = R(1-\cos\theta)\\
&a = g\tan\theta = \omega^2 r = \omega^2 R\sin\theta \qquad \cos\theta = \frac{g}{\omega^2 R}\\
&\omega\uparrow\quad H\uparrow\quad \omega\uparrow\ \text{上升．}
\end{align*}
% CHECK: 最后一行第三个符号手写同 ω（按推导可能是 θ↑）

% ---- block 076-04 (原稿第76页) kind=knowledge ----
\textbf{若不光滑．}
\begin{center}
\begin{tabular}{@{}ll@{}}
$a=g\tan\theta$ & 不受 $f$\\
$a>g\tan\theta$ & $f$ 向下\\
$a<g\tan\theta$ & $f$ 向上\\
\end{tabular}
\end{center}


\section{竖直面圆周·系统牛二}
% ---- block 087-08 (原稿第87页) kind=problem ----
\begin{problem}
\notefig[0.25]{figures/p087_g.png}
在最高点恰好对地面无压力，求在最低点压力
\begin{solution}
\[ (M+m)g=ma_{\text{上}}=m\,\frac{V_{\text{上}}^2}{R} \]
\[ N-(M+m)g=\frac{mV_{\text{下}}^2}{R}\qquad \text{在圆周必用动能定理} \]
\[ \frac12mV_{\text{下}}^2-\frac12mV_{\text{上}}^2=mg\,2R \]
\[ \therefore N=(2M+6m)g \]
\end{solution}
\end{problem}


\section{绳拉直问题}
% ---- block 088-03 (原稿第88页) kind=method ----
\noindent \textbf{绳拉直问题}\quad 沿绳方向速度消失
% “绳拉直问题”原稿被椭圆圈住
\notefig[0.45]{figures/p088_d.png}
\[ \begin{cases} V^2=2gL\sin\theta\\ mg(L-L\sin\theta)=\dfrac12mV'^2-\dfrac12m(V\cos\theta)^2\end{cases} \]
% CHECK: 第一式 “2gL sinθ” 中 g 与 L 之间有重叠笔迹（疑有一个 2：V^2=2g·2L sinθ）；图内标注为 “V sinθ 失去”“只剩 V cosθ”


\section{水平面圆周·质心与滑动临界}
% ---- block 089-02 (原稿第89页) kind=problem ----
\begin{problem}
\notefig[0.25]{figures/p089_a.png}
\begin{solution}
质心 $=\dfrac{2rm+mr}{2m}=\dfrac32r$ 距轴 $\dfrac12r$\qquad $2m\omega^2\dfrac r2>2\mu mg$ 时 向左滑。

若质心距轴为0 则不会滑动\qquad $2m\omega^2\cdot0=0$。
\end{solution}
\end{problem}


\section{平抛·圆轨道出口}
% ---- block 090-03 (原稿第90页) kind=problem ----
\begin{problem}
\notefig[0.35]{figures/p090_c.png}
求飞得最远的 $R$
\begin{solution}
\[ \begin{cases} S=V_0t\\ 2R=\dfrac12gt^2\\ \dfrac12mV^2-\dfrac12mV_0^2=-mg\,2R\end{cases}\qquad S=\sqrt{4R\left(\dfrac{V_0^2}{g}-4R\right)} \]
% CHECK: 第一式手写形似 S=V_0 t（按物理应为出口速度 V：S=Vt）
均值不等式\quad $R=\dfrac{V_0^2}{8g}$ 时取等
\end{solution}
\end{problem}


\section{斜抛·最大射程}
% ---- block 090-05 (原稿第90页) kind=derivation ----
\noindent \textbf{斜抛最大射程}\qquad 法1
\notefig[0.3]{figures/p090_e.png}
\[ \begin{cases} S=V_0\cos\theta\,t\\ -h=V_0\sin\theta\,t-\dfrac12gt^2\end{cases} \]
\[ -h=V_0\sin\theta\,\frac{S}{V_0\cos\theta}-\frac12g\,\frac{S^2}{(V_0\cos\theta)^2}\qquad\text{由}\ \frac{1}{\cos^2\theta}=1+\tan^2\theta \]
\[ \therefore\ \left(\frac{S}{V_0}\right)^2(1+\tan^2\theta)-\frac{2S}{g}\tan\theta-\frac{2h}{g}=0 \]
\[ \tan^2\theta-\frac{2SV_0^2}{gS^2}\tan\theta+\left(1-\frac{2hV_0^2}{gS^2}\right)=0 \]
\[ \Delta\ge0\ \ \text{故}\ \ S\le\sqrt{\frac{V_0^4}{g^2}+\frac{2V_0^2h}{g}}=\frac{V_0\sqrt{V_0^2+2gh}}{g}=\frac{V_0V_t}{g}\qquad\text{最大射程} \]
\noindent 最大高度\quad $L=\dfrac{(V_0\sin\theta)^2}{2g}$

% ---- block 090-06 (原稿第90页) kind=derivation ----
\noindent 法2
\notefig[0.25]{figures/p090_f.png}
\begin{align*}
V_t&=\sqrt{V_0^2+2gh}\\
x&=V_0t\cos\alpha\\
mgt&=m\Delta V\\
\Delta V&=V_0\sin\alpha+V_t\sin\beta\qquad V_0\cos\alpha=V_t\cos\beta
\end{align*}
\[ \therefore\ x=\frac{V_0V_t}{g}\sin(\alpha+\beta) \]
\[ \alpha+\beta=\frac{\pi}{2}\ \text{时}\ x_{\max}=\frac{V_0V_t}{g}\qquad \tan\alpha=\frac{V_0}{V_t} \]


\section{关联速度·绳端速度}
% ---- block 091-04 (原稿第91页) kind=derivation ----
\notefig[0.25]{figures/p091_c.png}

$(h^{2}+x^{2}=l^{2})'$\quad 隐函数求导
\[
2x\frac{\dd x}{\dd t}=2l\frac{\dd l}{\dd t}\qquad \therefore\ v=v_{\text{物}}\frac{x}{l}=v_{\text{物}}\cos\theta
\]
\[
v_{\text{物}}=\frac{v}{\cos\theta}=\frac{\sqrt{x^{2}+h^{2}}}{x}\,v
\qquad
a_{\text{物}}=\frac{\dd v_{\text{物}}}{\dd t}=\left(\frac{\dfrac{1}{2}\cdot\dfrac{2x^{2}v_{\text{物}}}{\sqrt{x^{2}+h^{2}}}-\sqrt{x^{2}+h^{2}}\,v_{\text{物}}}{x^{2}}\right)v
\]
% CHECK: a物 式中分子第一项前的 1/2 手写不太清楚


\section{竖直面匀速圆周运动·洗衣机}
% ---- block 099-04 (原稿第99页) kind=problem ----
④
\notefig[0.25]{figures/p099_e.png}

要衣物\unclear{贴}内壁则，$\omega\ge\sqrt{\dfrac{g}{R}}$

匀速圆周运动

洗衣机


\section{平抛}
% ---- block 114-01 (原稿第114页) kind=problem ----
% 本题只有标题“交弧用轨迹方程”、图和“求 Ek min”，没有写出完整题干
\begin{problem}
\textbf{交\unclear{弧}用轨迹方程}

\notefig[0.3]{figures/p114_a.png}

求 $E_{k\min}$
\begin{solution}
\begin{align*}
y&=-\frac{g}{2v_0^2}x^2+2h=\frac{x^2}{2h}\\
\therefore\ y&=\frac{2v_0^2h}{gh+v_0^2}
\end{align*}
先表达 $x^2=\dfrac{4hv_0^2h}{gh+v_0^2}=\dfrac{4v_0^2h^2}{gh+v_0^2}$

\struck{$y=\unclear{\ldots}$}\quad $y=\dfrac{1}{2h}x^2\ \Rightarrow\ y$．
\begin{align*}
E_k&=mg(2h-y)+\frac{1}{2}mv_0^2\qquad\swarrow\\
&=\frac{2mg^2h^2}{gh+v_0^2}+\frac{1}{2}mv_0^2\\
&=\frac{2mg^2h^2}{gh+v_0^2}+\frac{1}{2}m(gh+v_0^2)-\frac{1}{2}mgh\\
&\ge2\sqrt{2mg^2h^2\cdot\frac{1}{2}m}-\frac{1}{2}mgh=\frac{3}{2}mgh
\end{align*}
\end{solution}
\end{problem}


\section{圆周运动·同轴转动}
% ---- block 117-07 (原稿第117页) kind=formula ----
\begin{align*}
v&=\omega r & v&\propto r\\
a&=\frac{v^2}{r}=\omega^2 r & a&\propto r\\
T&=\frac{2\pi}{\omega}\propto\frac{1}{\omega}\\
\omega&=\frac{2\pi}{T}=2\pi f & f&\propto\omega\\
F_n&=m\omega^2 r\propto m\text{、}r.
\end{align*}


\section{曲线运动·轨迹与合力方向}
% ---- block 118-01 (原稿第118页) kind=knowledge ----
% 页面右上角露出邻页(化学)内容“混合气体 ρ=1.429g/L … 22.4 … 6.72L”，不属于本页，已忽略
\textbf{曲线运动}

速度合力夹轨迹\quad 凹侧方向偏向力

切线$v$\quad 凹侧力\quad 锐加钝减


\section{竖直面圆周·平抛临界}
% ---- block 195-03 (原稿第195页) kind=problem ----
% 原稿此块（及下方一行）被一条自右上向左下、末端回钩的大弧线划过，疑似划去，仍照录
\notefig[0.4]{figures/p195_b.png}

$R=0.4\unit{m}$

经$e$点$v$越小，下落高度越大，所用时间越长

物体刚好过$e$时速度最小

\begin{align*}
v_e&=\sqrt{gR}=2\unit{m/s}\\
x&=v_e t\\
y&=\tfrac{1}{2}gt^2
\end{align*}

\[
\frac{R-\dfrac{y-R}{\cos\theta}\ \star}{x-(y-R)\tan\theta}=\sin\theta
\]
% 分子右侧有一个小星号标记；分母括号内手写形似“4”，CHECK: 按几何关系应为 $y$

$\therefore\ t=0.3\unit{s}$

